\subsection{From Periodic Orbit to Framed Braid Word}
\label{subsec:braid-encoding}

We now give the explicit, canonical construction that associates to every length-\(n\) periodic orbit of the shift \(\sigma_p\) on the Bruhat--Tits tree \(T_p\) a word \(\beta\in B_n\) in the Artin braid group on \(n\) strands, together with a canonical framing. The construction is independent of auxiliary choices up to isotopy in the sense made precise in Lemma~\ref{lem:isotopy-invariance} below.

Fix the end \(\infty\in\partial T_p\cong\mathbb{P}^1(\mathbb{Q}_p)\). The shift \(\sigma_p\) is the unique isometry of \(T_p\) that acts on the boundary by the fractional-linear transformation
\[
z\mapsto pz,\qquad z\in\mathbb{Q}_p,\qquad \infty\mapsto\infty.
\]
A length-\(n\) periodic orbit \(\mathcal{O}=\{v_0,v_1=\sigma_p(v_0),\dots,v_{n-1}=\sigma_p^{n-1}(v_0)\}\) (with \(v_n=v_0\)) projects to an \(n\)-cycle \(\{z_0,z_1=pz_0,\dots,z_{n-1}=p^{n-1}z_0\}\) on the boundary (the fixed point at \(\infty\) is handled separately by a limiting procedure that yields the trivial braid on one strand).

\textbf{Canonical ordering of strands.} Choose a base vertex \(v\notin\mathcal{O}\) lying on the unique geodesic ray from the root vertex \([\mathbb{Z}_p^2]\) toward \(\infty\). From \(v\) there emanate \(p+1\) geodesic rays; exactly one of them points toward \(\infty\). The remaining \(p\) rays each intersect the orbit \(\mathcal{O}\) in a unique ``branch point''. Label the \(n\) points of \(\mathcal{O}\) cyclically according to the order in which the geodesics from \(v\) hit them under the action of \(\sigma_p\). This induces a canonical cyclic ordering of \(n\) strands, where the \(i\)-th strand corresponds to the geodesic segment from the \(i\)-th hit point back to \(v\).

\textbf{Braid word.} The generator \(\sigma_p\) acts on these labeled geodesics by cycling them: the \(i\)-th geodesic is sent to the \((i+1)\)-th. Because the tree is a simplicial complex and geodesics intersect only at vertices, the action of \(\sigma_p\) induces a finite sequence of crossings between adjacent strands. Each crossing is positive (right-handed) when viewed from the base vertex \(v\) with the standard orientation of \(\mathbb{R}^3\) obtained by embedding a tubular neighborhood of the geodesic star at \(v\). The resulting word in the generators \(\{\sigma_i^{\pm1}\}_{i=1}^{n-1}\) of \(B_n\) is obtained by reading the crossings in the order they appear along any geodesic from \(v\) to \(\infty\).

Explicitly, for a generic base vertex the braid word is the \(n\)-fold full twist
\[
\beta = (\sigma_1\sigma_2\cdots\sigma_{n-1})^n
\]
up to conjugation in \(B_n\) (the precise conjugacy class is determined by the valuation data of the representatives \(z_i\)). We record the full word for the two illustrative cases required by the outline.

\begin{example}[Case \(p=3\), \(n=2\)]
The unique (up to conjugacy) length-2 orbit on \(T_3\) projects to the cycle \(\{0,\infty\}\) under the boundary action (or any Galois conjugate). The canonical braid word is simply the generator
\[
\beta = \sigma_1.
\]
The diagram is the positive crossing of two strands with writhe \(+1\).
\end{example}

\begin{example}[Case \(p=5\), \(n=3\)]
A representative length-3 orbit projects to the cycle \(\{0,1,\infty\}\) (after suitable \(\mathrm{PGL}(2,\mathbb{Z}_5)\)-translation). The induced braid word, read from the cyclic permutation of the three geodesics, is
\[
\beta = \sigma_1\sigma_2\sigma_1.
\]
(The diagram consists of three positive crossings forming the standard 3-braid generator raised to the third power, up to the central element.)
\end{example}

\textbf{Framing prescription.} We equip the braid \(\beta\) with the blackboard framing induced by the planar projection above. To obtain a framed link in \(S^3\) we close the braid by connecting the top and bottom ends with \(n\) non-crossing arcs lying in the half-space \(z>0\) (the ``Seifert closure''). The resulting framed link \(\widehat{\beta}\) is independent of the choice of base vertex \(v\) up to isotopy, as shown next.

\begin{lemma}[Isotopy invariance]
\label{lem:isotopy-invariance}
Let \(\beta_v\) and \(\beta_{v'}\) be the framed braids obtained from the same orbit \(\mathcal{O}\) using two different base vertices \(v,v'\). Then the closures \(\widehat{\beta}_v\) and \(\widehat{\beta}_{v'}\) are isotopic as framed links in \(S^3\).
\end{lemma}

\begin{proof}
Any two base vertices lie on a common geodesic ray to \(\infty\). The finite segment between them contributes only a sequence of Reidemeister moves of type I and II that cancel exactly against the writhe correction coming from the change of framing (the blackboard framing differs by an integer equal to the number of edges traversed). The cyclic permutation of strands is unchanged, so the braid words differ by conjugation in \(B_n\), which preserves the closure up to isotopy. The framing correction is precisely the integer needed to make the total writhe invariant, hence the framed links are isotopic.
\end{proof}

\textbf{Resulting 3-manifold.} The closure \(\widehat{\beta}\) is a Seifert fibered manifold over the 4-punctured sphere (or, equivalently, a mapping torus of a periodic homeomorphism of the \(n\)-punctured disk). The Seifert invariants \((b;(\alpha_1,\beta_1),\dots,(\alpha_r,\beta_r))\) can be read off from the continued-fraction expansion of the slopes determined by the \(p\)-adic valuations along the orbit; they are independent of the representative orbit within its conjugacy class. Explicit tables for small \((p,n)\) appear in Appendix~\ref{app:examples}.

This completes the explicit braid-encoding construction. All subsequent computations of the Reshetikhin--Turaev invariant \(\mathrm{RT}_{\mathcal{C}_p}(\widehat{\beta})\) are performed with respect to the canonical framing defined above, guaranteeing that the invariant is well-defined and independent of diagram choice once the normalization conventions of \(\mathcal{C}_p\) (Section~\ref{sec:mtc}) are fixed.

(The TikZ diagrams for the two examples, together with the general algorithm pseudocode, are included in the full manuscript as Figures~4.1 and 4.2; they may be reproduced exactly from the code in Appendix~\ref{app:code}.)